\begin{abstract}
Scientific workflow runtimes must decide when a task is logically eligible,
where its inputs should be obtained, and when a worker should execute it.
Binding these decisions to a single admission event can leave computation
waiting for data coordination; advancing them independently can instead
create excessive staging and resource pressure. DataVine makes these three
decisions explicit. A logical scheduler dispatches eligible task descriptions,
a data controller owns intermediate identities and replica lifetimes, and
worker agents resolve inputs when execution opportunities approach. A bounded
worker queue is distinct from a feedback-controlled execution window. Physical
completion can release descendants before background backup, while requested
outputs have a separate durable-admission boundary. The implementation retains
individual task attempts and supports recall of queued, unstarted work.
Logical scheduling remains single-owner while Controller RPC handling and
durable-output persistence use independent data-plane concurrency pools.
We evaluate a complete ATLAS Open Data diphoton analysis, concurrency and
preparation controls and 2,048-task graphs across one
to eight distinct physical hosts. The analysis processes 36.56 million events
with exact agreement in all selection counts and histogram bins. Under matched
dependency preloading, DataVine elastic has a median paired completion-time
advantage of $\UpgradeAtlasElasticSpeedup{}\times$ over TaskVine. Fixed admission is
competitive with elastic admission; cold initialization and distributed
data-heavy chains expose substantial limitations. The artifact contains
\UpgradeTrials{} accepted trials in five randomized blocks per configuration
and \UpgradeRegressionPassed{} passing regression gates. The evidence supports
an explicit execution-boundary design, while limiting claims about adaptive
control and data-path scalability.\cite{codex}
\end{abstract}
\begin{IEEEkeywords}
scientific workflows, distributed runtime, intermediate data, admission control
\end{IEEEkeywords}
